This thesis presents the design and implementation of a credit risk-oriented business intelligence prototype that integrates data processing, multi-horizon prediction, experiment management, and analytical visualization within a social and solidarity economy institution in Ecuador.\\

The prototype features an integrated and automated workflow for data processing and predictive model updates, consolidates information into a datamart, and provides interactive visualizations, backed by a systematic experiment registry that ensures system traceability and observability.\\

The comparative evaluation among gradient boosting models, convolutional neural networks, and multilayer perceptrons shows that, under the studied configuration, the former exhibit superior discriminative capacity, although their performance degrades over extended time horizons.\\

The primary contribution lies in the technical integration of an end-to-end pipeline—processing, modeling, tracking, analytical storage, and visualization—that enables credit analysis and experiment observability. The prototype is functional in the evaluated environment; its operational validation and decision-impact assessment remain the responsibility of the institution under study.\\

\noindent\textbf{Keywords}:

credit risk; multi-horizon prediction; LightGBM; neural networks; business intelligence; MLOps; data mart; Apache Airflow.